% Options for packages loaded elsewhere
% Options for packages loaded elsewhere
\PassOptionsToPackage{unicode}{hyperref}
\PassOptionsToPackage{hyphens}{url}
\PassOptionsToPackage{dvipsnames,svgnames,x11names}{xcolor}
%
\documentclass[
  11pt,
]{article}
\usepackage{xcolor}
\usepackage[margin=1in]{geometry}
\usepackage{amsmath,amssymb}
\setcounter{secnumdepth}{-\maxdimen} % remove section numbering
\usepackage{iftex}
\ifPDFTeX
  \usepackage[T1]{fontenc}
  \usepackage[utf8]{inputenc}
  \usepackage{textcomp} % provide euro and other symbols
\else % if luatex or xetex
  \usepackage{unicode-math} % this also loads fontspec
  \defaultfontfeatures{Scale=MatchLowercase}
  \defaultfontfeatures[\rmfamily]{Ligatures=TeX,Scale=1}
\fi
\usepackage{lmodern}
\ifPDFTeX\else
  % xetex/luatex font selection
\fi
% Use upquote if available, for straight quotes in verbatim environments
\IfFileExists{upquote.sty}{\usepackage{upquote}}{}
\IfFileExists{microtype.sty}{% use microtype if available
  \usepackage[]{microtype}
  \UseMicrotypeSet[protrusion]{basicmath} % disable protrusion for tt fonts
}{}
\makeatletter
\@ifundefined{KOMAClassName}{% if non-KOMA class
  \IfFileExists{parskip.sty}{%
    \usepackage{parskip}
  }{% else
    \setlength{\parindent}{0pt}
    \setlength{\parskip}{6pt plus 2pt minus 1pt}}
}{% if KOMA class
  \KOMAoptions{parskip=half}}
\makeatother
% Make \paragraph and \subparagraph free-standing
\makeatletter
\ifx\paragraph\undefined\else
  \let\oldparagraph\paragraph
  \renewcommand{\paragraph}{
    \@ifstar
      \xxxParagraphStar
      \xxxParagraphNoStar
  }
  \newcommand{\xxxParagraphStar}[1]{\oldparagraph*{#1}\mbox{}}
  \newcommand{\xxxParagraphNoStar}[1]{\oldparagraph{#1}\mbox{}}
\fi
\ifx\subparagraph\undefined\else
  \let\oldsubparagraph\subparagraph
  \renewcommand{\subparagraph}{
    \@ifstar
      \xxxSubParagraphStar
      \xxxSubParagraphNoStar
  }
  \newcommand{\xxxSubParagraphStar}[1]{\oldsubparagraph*{#1}\mbox{}}
  \newcommand{\xxxSubParagraphNoStar}[1]{\oldsubparagraph{#1}\mbox{}}
\fi
\makeatother


\usepackage{longtable,booktabs,array}
\newcounter{none} % for unnumbered tables
\usepackage{calc} % for calculating minipage widths
% Correct order of tables after \paragraph or \subparagraph
\usepackage{etoolbox}
\makeatletter
\patchcmd\longtable{\par}{\if@noskipsec\mbox{}\fi\par}{}{}
\makeatother
% Allow footnotes in longtable head/foot
\IfFileExists{footnotehyper.sty}{\usepackage{footnotehyper}}{\usepackage{footnote}}
\makesavenoteenv{longtable}
\usepackage{graphicx}
\makeatletter
\newsavebox\pandoc@box
\newcommand*\pandocbounded[1]{% scales image to fit in text height/width
  \sbox\pandoc@box{#1}%
  \Gscale@div\@tempa{\textheight}{\dimexpr\ht\pandoc@box+\dp\pandoc@box\relax}%
  \Gscale@div\@tempb{\linewidth}{\wd\pandoc@box}%
  \ifdim\@tempb\p@<\@tempa\p@\let\@tempa\@tempb\fi% select the smaller of both
  \ifdim\@tempa\p@<\p@\scalebox{\@tempa}{\usebox\pandoc@box}%
  \else\usebox{\pandoc@box}%
  \fi%
}
% Set default figure placement to htbp
\def\fps@figure{htbp}
\makeatother





\setlength{\emergencystretch}{3em} % prevent overfull lines

\providecommand{\tightlist}{%
  \setlength{\itemsep}{0pt}\setlength{\parskip}{0pt}}



 


\usepackage{amsmath,amssymb,enumitem,booktabs}
\setlist{itemsep=2pt,topsep=3pt}
\usepackage{fancyhdr}\pagestyle{fancy}\fancyhf{}
\fancyhead[L]{\small DS701 Fall 2026}\fancyhead[R]{\small Homework 6}
\fancyfoot[C]{\small\thepage}
\renewcommand{\headrulewidth}{0.4pt}
\usepackage[breakable]{tcolorbox}
% --- answer-space macros (problems PDF only; no-ops in the solutions PDF) ---
% \answerspace{2in}  vertical room for handwritten work, sized from the solution
% \blank             a short fill-in rule; \blank[1.5in] for a wider one
% \answerpage        start the next problem on a fresh page (Gradescope page matching)
\newcommand{\answerspace}[1]{\par\vspace*{#1}}  % \vspace* so room is not dropped at a page break
\newcommand{\blank}[1][1.2in]{\rule{#1}{0.4pt}}
\newcommand{\answerpage}{\clearpage}
\newtcolorbox{solutionbox}{breakable,colback=blue!4,colframe=blue!45!black,title=Solution,fonttitle=\bfseries,left=4pt,right=4pt,top=3pt,bottom=3pt}
\makeatletter
\@ifpackageloaded{caption}{}{\usepackage{caption}}
\AtBeginDocument{%
\ifdefined\contentsname
  \renewcommand*\contentsname{Table of contents}
\else
  \newcommand\contentsname{Table of contents}
\fi
\ifdefined\listfigurename
  \renewcommand*\listfigurename{List of Figures}
\else
  \newcommand\listfigurename{List of Figures}
\fi
\ifdefined\listtablename
  \renewcommand*\listtablename{List of Tables}
\else
  \newcommand\listtablename{List of Tables}
\fi
\ifdefined\figurename
  \renewcommand*\figurename{Figure}
\else
  \newcommand\figurename{Figure}
\fi
\ifdefined\tablename
  \renewcommand*\tablename{Table}
\else
  \newcommand\tablename{Table}
\fi
}
\@ifpackageloaded{float}{}{\usepackage{float}}
\floatstyle{ruled}
\@ifundefined{c@chapter}{\newfloat{codelisting}{h}{lop}}{\newfloat{codelisting}{h}{lop}[chapter]}
\floatname{codelisting}{Listing}
\newcommand*\listoflistings{\listof{codelisting}{List of Listings}}
\makeatother
\makeatletter
\makeatother
\makeatletter
\@ifpackageloaded{caption}{}{\usepackage{caption}}
\@ifpackageloaded{subcaption}{}{\usepackage{subcaption}}
\makeatother
\usepackage{bookmark}
\IfFileExists{xurl.sty}{\usepackage{xurl}}{} % add URL line breaks if available
\urlstyle{same}
% fallback for those not using the hyperref driver hyperxmp:
\makeatletter
\@ifundefined{xmpquote}{\newcommand{\xmpquote}[1]{#1}}{}
\makeatother
\hypersetup{
  pdftitle={DS701 Homework 6: k-NN{,} the Curse and Naive Bayes},
  colorlinks=true,
  linkcolor={blue},
  filecolor={Maroon},
  citecolor={Blue},
  urlcolor={Blue},
  pdfcreator={LaTeX via pandoc}}


\title{DS701 Homework 6: k-NN, the Curse and Naive Bayes}
\usepackage{etoolbox}
\makeatletter
\providecommand{\subtitle}[1]{% add subtitle to \maketitle
  \apptocmd{\@title}{\par {\large #1 \par}}{}{}
}
\makeatother
\subtitle{Released Tue Oct 13 · due Mon Oct 19, 11:59 pm on Gradescope}
\author{}
\date{}
\begin{document}
\maketitle


\begin{center}\LARGE\bfseries\color{blue!45!black} SOLUTIONS\end{center}
\renewcommand{\answerspace}[1]{}
\renewcommand{\answerpage}{\par\vspace{6pt}}

\subsection{\texorpdfstring{Problem 1: \(k\)-NN by hand, and why scale
matters}{Problem 1: k-NN by hand, and why scale matters}}\label{problem-1-k-nn-by-hand-and-why-scale-matters}

\emph{Practices: the \(k\)-NN prediction rule, the effect of feature
units on Euclidean distance, and the bias--variance role of \(k\).}

Eight users of an app are described by \(x_1\) = hours of use per week
(0--10) and \(x_2\) = fraction of sessions rated ``satisfied'' (0--1).
Label A = renewed the subscription, B = cancelled.

{\def\LTcaptype{none} % do not increment counter
\begin{longtable}[]{@{}lllllllll@{}}
\toprule\noalign{}
point & P1 & P2 & P3 & P4 & P5 & P6 & P7 & P8 \\
\midrule\noalign{}
\endhead
\bottomrule\noalign{}
\endlastfoot
\(x_1\) & 4 & 8 & 6 & 0 & 2 & 7 & 10 & 1 \\
\(x_2\) & 0.2 & 0.9 & 0.7 & 0.2 & 0.0 & 0.1 & 1.0 & 0.5 \\
label & A & A & A & A & B & B & B & B \\
\end{longtable}
}

A new user has \(\mathbf{q} = (4,\ 0.1)\). Use Euclidean distance; since
only the \emph{ordering} of distances matters, you may work with squared
distances throughout.

\answerpage

\begin{enumerate}
\def\labelenumi{(\alph{enumi})}
\tightlist
\item
  Compute the squared distance from \(\mathbf{q}\) to each of the eight
  points. Give the \(k = 1\) and the \(k = 3\) predictions (majority
  vote).
\end{enumerate}

{\def\LTcaptype{none} % do not increment counter
\begin{longtable}[]{@{}
  >{\raggedright\arraybackslash}p{(\linewidth - 16\tabcolsep) * \real{0.1111}}
  >{\raggedright\arraybackslash}p{(\linewidth - 16\tabcolsep) * \real{0.1111}}
  >{\raggedright\arraybackslash}p{(\linewidth - 16\tabcolsep) * \real{0.1111}}
  >{\raggedright\arraybackslash}p{(\linewidth - 16\tabcolsep) * \real{0.1111}}
  >{\raggedright\arraybackslash}p{(\linewidth - 16\tabcolsep) * \real{0.1111}}
  >{\raggedright\arraybackslash}p{(\linewidth - 16\tabcolsep) * \real{0.1111}}
  >{\raggedright\arraybackslash}p{(\linewidth - 16\tabcolsep) * \real{0.1111}}
  >{\raggedright\arraybackslash}p{(\linewidth - 16\tabcolsep) * \real{0.1111}}
  >{\raggedright\arraybackslash}p{(\linewidth - 16\tabcolsep) * \real{0.1111}}@{}}
\toprule\noalign{}
\begin{minipage}[b]{\linewidth}\raggedright
point
\end{minipage} & \begin{minipage}[b]{\linewidth}\raggedright
P1
\end{minipage} & \begin{minipage}[b]{\linewidth}\raggedright
P2
\end{minipage} & \begin{minipage}[b]{\linewidth}\raggedright
P3
\end{minipage} & \begin{minipage}[b]{\linewidth}\raggedright
P4
\end{minipage} & \begin{minipage}[b]{\linewidth}\raggedright
P5
\end{minipage} & \begin{minipage}[b]{\linewidth}\raggedright
P6
\end{minipage} & \begin{minipage}[b]{\linewidth}\raggedright
P7
\end{minipage} & \begin{minipage}[b]{\linewidth}\raggedright
P8
\end{minipage} \\
\midrule\noalign{}
\endhead
\bottomrule\noalign{}
\endlastfoot
\(d^2\) & \blank[0.45in] & \blank[0.45in] & \blank[0.45in] &
\blank[0.45in] & \blank[0.45in] & \blank[0.45in] & \blank[0.45in] &
\blank[0.45in] \\
\end{longtable}
}

\(k = 1\) prediction: \blank[0.8in]\quad \(k = 3\) prediction:
\blank[0.8in]

\answerspace{1.5in}

\begin{enumerate}
\def\labelenumi{(\alph{enumi})}
\setcounter{enumi}{1}
\tightlist
\item
  A colleague points out that \(x_2\) is a fraction while \(x_1\) is in
  hours, and re-expresses \(x_2\) as a score out of 10 (multiply the
  \(x_2\) column and \(q_2\) by 10). Recompute the squared distances,
  and give the new \(k = 1\) and \(k = 3\) predictions. Which prediction
  changed, and why exactly --- which point(s) entered or left the
  neighbourhood? In one sentence, what does this say about running
  \(k\)-NN on raw features?
\end{enumerate}

New \(k = 1\) prediction: \blank[0.8in]\quad new \(k = 3\) prediction:
\blank[0.8in]

\answerspace{2.8in}
\answerpage

\begin{enumerate}
\def\labelenumi{(\alph{enumi})}
\setcounter{enumi}{2}
\tightlist
\item
  What does \(k = 8\) predict for \emph{any} query in this dataset, and
  why? Using the vocabulary of Lecture 07 (bias/variance), explain what
  moving \(k\) from 1 toward 8 does to the decision boundary and to the
  two kinds of error, and say how \(k\) should be chosen in practice.
\end{enumerate}

\(k = 8\) predicts: \blank[1.2in]

\answerspace{2in}

\begin{solutionbox}

\begin{enumerate}
\def\labelenumi{(\alph{enumi})}
\item
  \(d^2(\mathbf{q}, P_i) = (x_1-4)^2 + (x_2-0.1)^2\): P1:
  \(0 + 0.01 = 0.01\); P2: \(16 + 0.64 = 16.64\); P3:
  \(4 + 0.36 = 4.36\); P4: \(16 + 0.01 = 16.01\); P5:
  \(4 + 0.01 = 4.01\); P6: \(9 + 0 = 9.00\); P7: \(36 + 0.81 = 36.81\);
  P8: \(9 + 0.16 = 9.16\). Sorted: P1 (A), P5 (B), P3 (A), P6 (B), P8
  (B), P4 (A), P2 (A), P7 (B). \(k = 1\): P1 \(\Rightarrow\) \textbf{A}.
  \(k = 3\): \{P1, P5, P3\} = A, B, A \(\Rightarrow\) \textbf{A}.
\item
  Now \(x_2 \in \{2, 9, 7, 2, 0, 1, 10, 5\}\) and \(q_2 = 1\);
  \(d^2 = (x_1-4)^2 + (x_2-1)^2\): P1: \(0 + 1 = 1\); P2:
  \(16 + 64 = 80\); P3: \(4 + 36 = 40\); P4: \(16 + 1 = 17\); P5:
  \(4 + 1 = 5\); P6: \(9 + 0 = 9\); P7: \(36 + 81 = 117\); P8:
  \(9 + 16 = 25\). Sorted: P1 (A), P5 (B), P6 (B), P4 (A), P8 (B), P3
  (A), P2 (A), P7 (B). \(k = 1\): still P1 \(\Rightarrow\) \textbf{A}.
  \(k = 3\): \{P1, P5, P6\} = A, B, B \(\Rightarrow\) \textbf{B} ---
  flipped. Before rescaling, \(x_2\) differences were at most \(1\)
  against \(x_1\) differences of up to \(10\), so \(x_2\) was
  effectively ignored and P3 (\(x_1 = 6\), but a very different
  satisfaction, \(0.7\) vs \(0.1\)) counted as close. After rescaling,
  P3's satisfaction gap costs \(36\) and it drops to sixth; P6 (similar
  satisfaction, \(x_1 = 7\)) takes its place. \(k\)-NN's answer depends
  on the \emph{units} you happened to record --- so scale (standardize)
  features first; otherwise the largest-range feature silently decides.
  \emph{Common mistake:} comparing distances instead of squared
  distances is fine, but mixing the two within one ranking is not.
\item
  \(k = 8\) uses every point, so it always returns the majority class of
  the whole training set --- here a 4--4 tie (any tie-break rule; the
  point is that the query's location no longer matters). Small \(k\):
  jagged boundary that follows individual points, low bias, high
  variance (a different training sample moves the boundary a lot). Large
  \(k\): smooth boundary, low variance, high bias --- at \(k = N\) it
  ignores \(\mathbf{x}\) entirely. Choose \(k\) on held-out data
  (cross-validation over a grid of \(k\)), never by training error,
  which is trivially perfect at \(k = 1\).
\end{enumerate}

\end{solutionbox}

\answerpage

\subsection{Problem 2: The curse of dimensionality, in
numbers}\label{problem-2-the-curse-of-dimensionality-in-numbers}

\emph{Practices: quantifying ``nearest is meaningless'' with two short
calculations.}

Suppose \(N\) points are spread uniformly in the unit hypercube
\([0,1]^d\). To find the \(k\) nearest neighbours of a query, \(k\)-NN
needs a small box around the query that contains a fraction \(f = k/N\)
of the data.

\begin{enumerate}
\def\labelenumi{(\alph{enumi})}
\tightlist
\item
  Show that a hypercube of side length \(s\) inside \([0,1]^d\) captures
  a fraction \(s^d\) of the data, so the side length needed to capture
  fraction \(f\) is \(s = f^{1/d}\). (Two or three lines.)
\end{enumerate}

\answerspace{3.8in}

\begin{enumerate}
\def\labelenumi{(\alph{enumi})}
\setcounter{enumi}{1}
\tightlist
\item
  For \(f = 0.1\) compute \(s\) for \(d = 1, 2, 10, 100\). You may use
  \(0.1^{0.5} \approx 0.316\), \(0.1^{0.1} \approx 0.794\),
  \(0.1^{0.01} \approx 0.977\). In one or two sentences, what does the
  \(d = 100\) answer mean for the word ``neighbourhood'' in \(k\)-NN?
\end{enumerate}

{\def\LTcaptype{none} % do not increment counter
\begin{longtable}[]{@{}lllll@{}}
\toprule\noalign{}
\(d\) & 1 & 2 & 10 & 100 \\
\midrule\noalign{}
\endhead
\bottomrule\noalign{}
\endlastfoot
\(s\) & \blank[0.8in] & \blank[0.8in] & \blank[0.8in] & \blank[0.8in] \\
\end{longtable}
}

\answerspace{1.2in}
\answerpage

\begin{enumerate}
\def\labelenumi{(\alph{enumi})}
\setcounter{enumi}{2}
\tightlist
\item
  Lecture 09 showed the fraction of a unit ball's volume within
  \(\epsilon\) of its surface is \(1 - (1-\epsilon)^d\). For
  \(\epsilon = 0.1\) compute it at \(d = 2, 10, 50\) (use
  \(0.9^{10} \approx 0.349\), \(0.9^{50} \approx 0.005\)). Combine this
  with the fact that \(d_{\max}/d_{\min} \to 1\): why does ``the nearest
  neighbour'' stop carrying information as \(d\) grows?
\end{enumerate}

Fraction within \(\epsilon = 0.1\) of the surface:

\(d = 2\): \blank[0.8in]\quad \(d = 10\): \blank[0.8in]\quad \(d = 50\):
\blank[0.8in]

\answerspace{2.8in}

\begin{enumerate}
\def\labelenumi{(\alph{enumi})}
\setcounter{enumi}{3}
\tightlist
\item
  A teammate proposes to fix a poorly performing \(k\)-NN model on 500
  features by ``just collecting more data''. Using a
  10-bins-per-dimension argument (with \(10^6\) points, how many points
  per bin at \(d = 3\)? at \(d = 6\)? at \(d = 12\)?), say why that plan
  fails, and name the two remedies the lecture actually recommends.
\end{enumerate}

Points per bin --- \(d = 3\): \blank[0.8in]\quad \(d = 6\):
\blank[0.8in]\quad \(d = 12\): \blank[0.8in]

\answerspace{1.8in}

\begin{solutionbox}

\begin{enumerate}
\def\labelenumi{(\alph{enumi})}
\item
  A box of side \(s\) has volume \(s^d\); the unit cube has volume 1.
  Uniform points land in a region with probability equal to its volume,
  so the expected fraction of points in the box is \(s^d/1 = s^d\).
  Setting \(s^d = f\) gives \(s = f^{1/d}\).
\item
  \(d = 1\): \(0.1\); \(d = 2\): \(0.316\); \(d = 10\): \(0.794\);
  \(d = 100\): \(0.977\). In 100 dimensions, a box holding just \(10\%\)
  of the data must span \(97.7\%\) of the range of \emph{every}
  coordinate --- it is essentially the whole cube. The ``neighbourhood''
  is not local: the \(k\) neighbours are typical points, not points that
  resemble the query.
\item
  \(d = 2\): \(1 - 0.81 = 0.19\); \(d = 10\): \(1 - 0.349 = 0.651\);
  \(d = 50\): \(1 - 0.005 =
  0.995\). So in 50-D, \(99.5\%\) of everything ``within distance
  \(r\)'' of the query sits in the outer \(10\%\) shell --- at distance
  \(\approx r\). Every candidate neighbour is about equally far, which
  is exactly \(d_{\max}/d_{\min} \to 1\): the closest point is barely
  closer than the farthest, so which point is ``nearest'' is decided by
  noise and its label tells you little about the query.
\item
  \(10^6\) points, \(10^d\) bins: \(d = 3\): \(10^6/10^3 = 1000\) points
  per bin; \(d = 6\): \(1\) point per bin; \(d = 12\): \(10^{-6}\),
  i.e.~one point per million bins. Keeping the density fixed as \(d\)
  grows needs \(10^d\) points --- exponential in \(d\); at \(d = 500\)
  no dataset is remotely large enough, so ``more data'' cannot make the
  space dense. Remedies from the lecture: use a different metric less
  affected by concentration (e.g.~cosine distance), or reduce the
  dimension first (feature selection / dimensionality reduction such as
  PCA, coming in Lecture 13).
\end{enumerate}

\end{solutionbox}

\answerpage

\subsection{Problem 3: Naive Bayes on symptoms, and the SVM margin as
geometry}\label{problem-3-naive-bayes-on-symptoms-and-the-svm-margin-as-geometry}

\emph{Practices: Bayes-rule classification by counting, Laplace
smoothing, and the distance-to-a-hyperplane computation behind support
vector machines.}

\textbf{Part I --- Naive Bayes.} A clinic recorded three symptoms and a
lab-confirmed diagnosis for ten patients:

{\def\LTcaptype{none} % do not increment counter
\begin{longtable}[]{@{}lllll@{}}
\toprule\noalign{}
\# & Fever & Cough & Headache & Flu \\
\midrule\noalign{}
\endhead
\bottomrule\noalign{}
\endlastfoot
1 & Y & Y & Y & Y \\
2 & Y & Y & Y & Y \\
3 & Y & Y & Y & Y \\
4 & Y & Y & N & Y \\
5 & Y & Y & Y & N \\
6 & N & Y & Y & N \\
7 & N & Y & N & N \\
8 & N & N & N & N \\
9 & N & N & N & N \\
10 & N & N & N & N \\
\end{longtable}
}

A new patient presents with \textbf{Fever = Y, Cough = N, Headache = Y}.

\begin{enumerate}
\def\labelenumi{(\alph{enumi})}
\tightlist
\item
  Estimate the class priors and the six conditional probabilities Naive
  Bayes needs for this patient,
  \(P(\text{symptom value} \mid \text{Flu})\) and
  \(P(\text{symptom value} \mid \text{no Flu})\), by counting. Compute
  the Naive Bayes score \(P(C)\prod_j P(x_j \mid C)\) for each class and
  give the prediction. What single number decided it?
\end{enumerate}

{\def\LTcaptype{none} % do not increment counter
\begin{longtable}[]{@{}
  >{\raggedright\arraybackslash}p{(\linewidth - 10\tabcolsep) * \real{0.1667}}
  >{\raggedright\arraybackslash}p{(\linewidth - 10\tabcolsep) * \real{0.1667}}
  >{\raggedright\arraybackslash}p{(\linewidth - 10\tabcolsep) * \real{0.1667}}
  >{\raggedright\arraybackslash}p{(\linewidth - 10\tabcolsep) * \real{0.1667}}
  >{\raggedright\arraybackslash}p{(\linewidth - 10\tabcolsep) * \real{0.1667}}
  >{\raggedright\arraybackslash}p{(\linewidth - 10\tabcolsep) * \real{0.1667}}@{}}
\toprule\noalign{}
\begin{minipage}[b]{\linewidth}\raggedright
class \(C\)
\end{minipage} & \begin{minipage}[b]{\linewidth}\raggedright
\(P(C)\)
\end{minipage} & \begin{minipage}[b]{\linewidth}\raggedright
\(P(F{=}Y \mid C)\)
\end{minipage} & \begin{minipage}[b]{\linewidth}\raggedright
\(P(C{=}N \mid C)\)
\end{minipage} & \begin{minipage}[b]{\linewidth}\raggedright
\(P(H{=}Y \mid C)\)
\end{minipage} & \begin{minipage}[b]{\linewidth}\raggedright
score
\end{minipage} \\
\midrule\noalign{}
\endhead
\bottomrule\noalign{}
\endlastfoot
Flu & \blank[0.7in] & \blank[0.7in] & \blank[0.7in] & \blank[0.7in] &
\blank[0.7in] \\
no Flu & \blank[0.7in] & \blank[0.7in] & \blank[0.7in] & \blank[0.7in] &
\blank[0.7in] \\
\end{longtable}
}

Prediction: \blank[0.8in]

\answerspace{1.5in}
\newpage

\begin{enumerate}
\def\labelenumi{(\alph{enumi})}
\setcounter{enumi}{1}
\tightlist
\item
  Redo the scores with Laplace smoothing \(\alpha = 1\): each
  conditional becomes
  \(\dfrac{\#\{x_j = v \text{ and class } C\} + \alpha}{\#\{\text{class } C\} + \alpha V_j}\),
  where \(V_j = 2\) is the number of values symptom \(j\) can take
  (leave the priors unsmoothed). Give the smoothed prediction and the
  smoothed posterior \(P(\text{Flu} \mid \mathbf{x})\) = (Flu score) /
  (sum of both scores). Did smoothing change the prediction, the
  confidence, or both? Round to 3 decimals.
\end{enumerate}

Smoothed prediction:
\blank[0.8in]\quad \(P(\text{Flu} \mid \mathbf{x})\): \blank[0.8in]

\answerspace{2.4in}

\begin{enumerate}
\def\labelenumi{(\alph{enumi})}
\setcounter{enumi}{2}
\tightlist
\item
  Suppose the clinic added a fourth column, ``Chills'', that happens to
  be identical to Fever in every row. Say what Naive Bayes would do to
  the fever evidence, why that violates the ``naive'' assumption, and
  --- in one or two sentences --- why the classifier often still picks
  the right class despite such violations.
\end{enumerate}

\answerspace{1.8in}

\begin{solutionbox}

\begin{enumerate}
\def\labelenumi{(\alph{enumi})}
\item
  Priors: \(P(\text{Flu}) = 4/10 = 0.4\), \(P(\text{no}) = 6/10 = 0.6\).
  Flu rows 1--4: \(P(F{=}Y \mid \text{Flu}) = 4/4 = 1\),
  \(P(C{=}N \mid \text{Flu}) = 0/4 = 0\),
  \(P(H{=}Y \mid \text{Flu}) = 3/4\). No-flu rows 5--10:
  \(P(F{=}Y \mid \text{no}) = 1/6\),
  \(P(C{=}N \mid \text{no}) = 3/6 = 1/2\),
  \(P(H{=}Y \mid \text{no}) = 2/6 = 1/3\). Scores: Flu:
  \(0.4 \cdot 1 \cdot 0 \cdot 0.75 = 0\); no:
  \(0.6 \cdot \tfrac16 \cdot \tfrac12 \cdot \tfrac13 = 0.0167\).
  Prediction: \textbf{no flu} --- decided entirely by the single zero:
  no training flu patient lacked a cough, so ``no cough'' vetoes flu
  regardless of fever and headache.
\item
  Smoothed conditionals (\(+1\) over \(+2\)): Flu:
  \(P(F{=}Y) = 5/6 = 0.833\), \(P(C{=}N) = 1/6 = 0.167\),
  \(P(H{=}Y) = 4/6 = 0.667\); no: \(P(F{=}Y) = 2/8 = 0.25\),
  \(P(C{=}N) = 4/8 = 0.5\), \(P(H{=}Y) = 3/8 = 0.375\). Scores: Flu:
  \(0.4 \cdot 0.833 \cdot 0.167 \cdot 0.667 = 0.037\); no:
  \(0.6 \cdot 0.25 \cdot 0.5 \cdot 0.375 = 0.028\). Prediction flips to
  \textbf{Flu}; posterior
  \(P(\text{Flu} \mid \mathbf{x}) = 0.037/(0.037+0.028) = 0.568\). Both
  changed: the veto is gone (an unseen combination is now merely
  unlikely, not impossible), and the fever evidence --- \(0.833\) vs
  \(0.25\), a factor above 3 --- is allowed to weigh in. Note the modest
  posterior: with four flu rows the estimates are crude, and smoothing
  says so.
\item
  Naive Bayes multiplies in \(P(\text{Chills} \mid C)\) as an
  independent factor, which equals \(P(\text{Fever} \mid C)\) --- it
  counts the fever evidence \emph{twice} (the fever likelihood ratio
  gets squared). That is exactly what conditional independence forbids:
  given the class, knowing Fever tells you Chills with certainty. It
  often does not matter because classification only needs the \emph{arg
  max} to be right; double-counting distorts the probabilities
  (over-confidence) but tends to push the score further in the direction
  the evidence already pointed, so the winner is usually unchanged.
\end{enumerate}

\end{solutionbox}




\end{document}
